% \begin{figure*}[!t]  
%     \centering  
%     \includegraphics[width=17cm,height=3.13cm]{figs/source_data_statistics.png} % Replace 'example-image' with your image file name  
%     \caption{The statistical analysis of the source data.
%     This analysis elucidates the distribution of various video types, including films, television shows, and documentaries with various time durations.
%     Additionally, we examine the categorization of these videos, their resolutions, and the keywords associated with textual captions.
%     \textcolor{red}{@Liwei}}
%     \label{fig:blank_image}  
% \end{figure*} 

\begin{figure*}[!t]
    \centering
    \begin{subfigure}{0.23\linewidth}
        \centering
        \includegraphics[width=1.0\textwidth]{figs/source_data_category.png} 
        \caption{Video categories}  
        \label{fig:source_data_category}  
    \end{subfigure}%
    \hfill
    \begin{subfigure}{0.23\linewidth}
        \centering
        \includegraphics[width=1.0\textwidth]{figs/source_data_duration.png}  
        \caption{Video duration}  
        \label{fig:source_data_duration}  
    \end{subfigure}%
    \hfill
    \begin{subfigure}{0.23\linewidth}
        \centering
        \includegraphics[width=1.0\textwidth]{figs/source_data_resolution.png}  
        \caption{Video resolution}  
        \label{fig:source_data_resolution}  
    \end{subfigure}  
    \begin{subfigure}{0.23\linewidth}
        \centering
        \includegraphics[width=1.0\textwidth]{figs/source_data_keywords.png}  
        \caption{Caption keywords}  
        \label{fig:source_data_keywords}  
    \end{subfigure}%
    % \vspace{-2mm}
    \caption{The statistical analysis of the source data.
    This analysis elucidates the distribution of various video types, including films, television shows, and documentaries of various time durations.
    Additionally, we examine the categorization of these videos, their resolutions, and the keywords associated with textual captions.}
    \label{fig:source_data}
    % \vspace{-2mm}
\end{figure*}

\begin{table*}[!t]
\centering
\resizebox{0.95\textwidth}{!}{
\begin{tabular}{@{}ccccccccc@{}}
\toprule
 Dataset name & Year&Domain & \# Videos &Total length (hours)&Caption type&Motion type& Resolution \\ 
\midrule
 WebVid-10M\cite{bain2021frozen}&2021&Open&10M&52K&Short&Text&360P\\
  Panda-70M\cite{chen2024panda}&2024&Open&70M&167K&Short&Text&720P\\
  OpenVid-1M\cite{nan2024openvid}&2024&Open&1M&2K&Long&Text&512P\\
  Koala-36M\cite{wang2024koala}&2024&Open&36M&172K&Long&Text&720P\\
  \midrule
  UCF-101\cite{soomro2012ucf101}& 2012&Human action & 13.3K & 26.7&Short&Text&240P\\
  NTU RGB+D\cite{shahroudy2016ntu}&2014&Human action&114K&37&-&3D pose, depth&1080P\\
  MSP-Avatar\cite{sadoughi2015msp}&2015&Human action&74&3&-&Speech audio,  landmark, pose&1080P\\
  ActivityNet\cite{caba2015activitynet}  & 2017&Human action & 100K & 849 & Short& Text& - \\
  TikTok-v4 \cite{chang2023magicdance} & 2023 &Human dance& 350 & 1& -&Skeleton & -  \\
  \midrule
Ours&2024&Human&52.3M&70.6K&Short, long, structured&Text, skeleton pose, speech audio&720P\\
Ours (filtered)&2024&Human&13.2M&16.7K&Short, long, structured&Text, skeleton pose, speech audio&720P\\
\bottomrule
\end{tabular}
}
% \vspace{-2mm}
\caption{The comparative analysis of our dataset against previous general and human video datasets. 
We enhance the textual captions by incorporating short, long, and structured formats that reflect human characteristics. 
Additionally, we integrate skeleton sequences derived from DWPose~\cite{yang2023effective} and corresponding speech audio filtered through SyncNet~\cite{raina2022syncnet} to enrich the dataset with contextual human motion data.}
% \vspace{-5.5mm}
\label{tab:datasets}
\end{table*}

\section{Related Work}
\noindent\textbf{Diffusion Video Generation.}
Video generation primarily utilizes two methodologies: diffusion-based and autoregressive methods. Diffusion-based methods can be further categorized into U-Net-based and Diffusion Transformer (DiT)-based architectures, depending on the design of the denoising network.
U-Net-based diffusion models have seen significant advancements through seminal works such as Video Diffusion Models (VDMs)~\cite{ho2022video}, Imagen Video~\cite{ho2022imagen}, and Align Your Latents~\cite{blattmann2023align}. These approaches extend traditional image diffusion architectures to accommodate temporal data by integrating temporal dimensions into the latent space or by jointly training on image and video datasets. Recent open-source models, including Stable Video Diffusion~\cite{blattmann2023stable} and ModelScopeT2V~\cite{wang2023modelscope} have further improved accessibility and performance, enabling high-resolution video synthesis with enhanced temporal coherence.
In contrast, DiT-based methods, exemplified by models such as Sora~\cite{liu2024sora}, Open-Sora~\cite{opensora}, Vidu~\cite{bao2024vidu}, and CogVideoX~\cite{yang2024cogvideox}, leverage transformer architectures that operate on spatiotemporal patches of video latent codes. This strategy facilitates the generation of longer, higher-resolution videos with improved coherence and dynamism by effectively capturing long-range dependencies in both spatial and temporal dimensions. The emergence of diffusion transformers and autoregressive models highlights the increasing demand for large-scale, high-quality data, which this paper addresses within the human domain.

\begin{figure*}[htbp]
    \centering
    \includegraphics[width=0.93\textwidth]{figs/datapipeline1.png} 
    % \vspace{-2.5mm}
    \caption{
The data processing pipeline.
The inputs are 105K hours of raw data from films, television shows, and documentaries, 
and the outputs are filtered high-quality videos that include textual captions—both short and long, as well as structured captions containing human information—and specific motion conditions related to individuals, such as skeleton sequences and speech audio. 
This pipeline consists of four key steps: video preprocessing, which involves basic decoding, cropping, and segmentation of the video; 
and video quality filtering, which assesses various quality metrics including luminance, blur, aesthetics, motion and technical quality. 
Then the human skeleton and speech audio are extracted from the video clips. 
An initial textual caption are generated with MiniCPM and CogVLM, voting by BLIP2 and reorgnized by the classic Llama model to obtain textual captions of different types.
Furthermore, the pipeline incorporates an advanced human quality filtering stage that aligns textual captions with the appearance, expressions, and pose movements of individuals, promoting fine-grained alignment between the textual information and the visual characteristics of the subjects.}
\label{fig:data_pipeline}
% \vspace{-mm}
\end{figure*}

\noindent\textbf{General Video Dataset.}
High-quality video-language datasets are crucial for advancing general video generation tasks. 
Notable large-scale datasets, such as HD-VILA-100M~\cite{xue2022advancing} and HowTo100M~\cite{miech2019howto100m}, utilize automatic speech recognition (ASR) for annotation; however, the captions generated often fail to adequately represent the video content, thereby limiting their utility for training purposes. 
WebVid-10M~\cite{bain2021frozen} stands out as a pioneering open-scenario text-to-video dataset, comprising 10.7 million text-video pairs with a cumulative duration of 52,000 hours. 
Panda-70M~\cite{chen2024panda} assembles 70 million high-resolution video samples characterized by strong semantic coherence. 
InternVid~\cite{wang2023internvid} leverages large language models to autonomously construct a video-text dataset, resulting in the generation of 234 million video clips accompanied by text descriptions. 
OpenVid-1M~\cite{nan2024openvid} provides expressive captions and curates 433,000 1080P videos, and its high-quality subset, OpenVidHD-0.4M, enhances high-resolution video generation.
Koala-36M~\cite{wang2024koala} includes 36 million video clips at a resolution of 720P, offering structured captions that average 200 words per segmented video clip, thereby improving the alignment between text and video content. 
Despite the advancements represented by these datasets, they exhibit a relative lack of diversity in terms of identity and motion within the human category, as well as deficiencies in high-quality textual captions, skeleton representations, and speech audio as motion conditions.

% \noindent\textbf{Human Video Dataset.}
% To support human image animation tasks, diverse and large-scale human-centric video datasets are essential. Real-world datasets collected from the Internet, such as the TikTok~\cite{jafarian2021learning}, contains over 300 dance video sequences shared on a social media, totaling more than 100,000 images. To expand dataset sizes cost-effectively, several synthetic datasets have been introduced. For example, SURREAL~\cite{varol2017learning} is a large-scale dataset featuring synthetically-generated but realistic images of people rendered from 3D sequences of human motion capture data. HSPACE~\cite{bazavan2021hspace} generates animations by fitting a human body model, offering diversity in characters, motions, and scenes. SynBody~\cite{yang2023synbody} is a large-scale synthetic dataset for human perception and modeling, sampling 10,000 animatable subjects. BEDLAM~\cite{black2023bedlam} renders people in realistic scenes and uses physics simulations to achieve realistic clothing effects. HumanVid~\cite{wang2024humanvid} introduces a synthetic dataset combined with a high-quality real-world video dataset, specifically designed for human image animation.
% While these datasets excel in specific tasks, they are not well-suited for pre-training or fine-tuning general video generative models in the human domain.
\noindent\textbf{Human Video Datasets.}
Diverse and large-scale human-centric video datasets are essential for advancing human image animation tasks. Real-world datasets collected from the Internet, such as TikTok~\cite{jafarian2021learning}, contain over 300 dance video sequences shared on social media platforms, totaling more than 100,000 images. UCF101~\cite{soomro2012ucf101} is a widely used action recognition dataset comprising realistic action videos collected from YouTube, organized into 101 action categories. NTU RGB+D~\cite{shahroudy2016ntu} is a large-scale dataset specifically designed for human action recognition, consisting of 56,880 samples spanning 60 action classes. Similarly, ActivityNet~\cite{caba2015activitynet} offers samples from 203 activity classes, totaling of 849 hours of video content. BEDLAM~\cite{black2023bedlam} generates realistic scenes of people using physics simulations to create accurate clothing effects, while HumanVid~\cite{wang2024humanvid} integrates synthetic and high-quality real-world video datasets tailored specifically for human image animation. Despite the utility of these datasets for specific tasks, they remain inadequate for pre-training or fine-tuning general video generative models in the human domain due to limitations in diversity, scale, and the specific motion conditions required for such models.